\newif\ifglobal

%%%%%%%%%%%%%%%%%%%%%%%
%%% Compile to view %%%
%%%%%%%%%%%%%%%%%%%%%%%

%%%%%%%%%%%%%%%%%%%%%%%%%%%%%%%%%%%%%%%%%%%%%%%%%%%
%%% Readme:                                     %%%
%%% https://www.overleaf.com/read/bwdjvfjksvjy/ %%%
%%%%%%%%%%%%%%%%%%%%%%%%%%%%%%%%%%%%%%%%%%%%%%%%%%%

% >>> M?: marks a module setting number or important notice

% >>> M4: Set {./relative-path} to external files for this module <<<
\def \modulefiles {./62-TMPSEN}
% To add an external file, use:
% (Replace '---' with actual filename)
% '\input{\modulefiles/tables/---.tex}' for tables
% '\includegraphics{\modulefiles/figures/---}' for figures
% '\subfile{\modulefiles/---}' for register files

% >>> M1: This module is in global project? <<<
% 'YES' - uncomment; 'NO' - comment out
\globaltrue

\ifglobal
    \documentclass[main\_\_EN.tex]{subfiles}

\else
    \documentclass[openany, 10pt]{book}

    % >>> M2: In ./chip-spec-settings/preamble-trm-module, also find
    %         settings for visibility of todo notes, watermark <<<
    \usepackage{preamble-trm-module}

    % >>> M3: Temporary labels in English,
    %         you can add your own labels there <<<
    \myexternaldocument{temp-labels-en}

\fi %\ifglobal


\setcounter{secnumdepth}{4}




% >>> M5: If this module needs any updates in
% './chip-spec-settings/preamble-trm-module.sty'
% (include additional package, etc.), contact your module PM
% or create the file './chip-spec-settings/changelog.tex' make the
% necessary updates yourself and document them in this file <<<

\raggedright


\begin{document}


% Sets language into which pre-defined bits of text are translated
\selectlanguage{english}

% Do not delete this block
\ifglobal
% Add the following front matter if
% this module is outside of global project
\else
    \InputIfFileExists{readme}{}{}
    {\let\clearpage\relax \listoftodos}
    {\let\clearpage\relax \listofdonetodos}
    \tableofcontents
    %\listoftables
    %\listoffigures
\fi


%%%%%%%%%%%%%%%%%%%%%%%%%
%%% TRM Module Begins %%%
%%%%%%%%%%%%%%%%%%%%%%%%%

\hypertarget{tmpsen}{}
\chapter{Temperature Sensor (TSENS)}
\label{mod:tmpsen}

\section{Overview}

\chipname{} provides a temperature sensor for real-time monitoring of temperature changes within the chip. The sensor converts analog voltage to digital values and provides compensation for temperature offsets.

\section{Features}

The temperature sensor has the following features:

\begin{itemize}
    \item Software-triggered temperature measurement, which once triggered, the sensor continuously measures temperature. Software can read the data at any time.
    \item Hardware-triggered automatic temperature monitoring, supporting two wake-up modes
    \item Configurable temperature offset based on the application scenario for improved accuracy
    \item Configurable temperature measurement range
\item Support for Event Task Matrix (ETM)-related events and tasks
\end{itemize}

\section{Architecture}

Figure \ref{fig:tsens_struct} shows the internal structure of the temperature sensor.

\begin{figure}[H]
    \centering
    \includegraphics[width=0.8\textwidth]{\modulefiles/figures/tsens_struct_20221209}
    \caption{Temperature Sensor Architecture}
    \label{fig:tsens_struct}
\end{figure}

As Figure \ref{fig:tsens_struct} shows, the temperature sensor module contains the following major blocks:

\begin{itemize}
\item Tsens Ctrl \& Detect: temperature sensor
    \item HW Timer: triggers automatic temperature monitoring
    \item Tsens\_data Monitor: monitors whether the temperature is outside the threshold range
    \item sync\_module: Synchronization block between APB clock domain and temperature sensor clock domain
\end{itemize}

\section{Functional Description}

\subsection{Temperature Sensor Power Up}

The temperature sensor can be powered up by setting the \hyperref[fielddesc:TSENSPOWERUP]{TSENS\_POWER\_UP} register.

\subsection{Temperature Sensor Clock}

The temperature sensor has only one clock source: LP\_PERI\_CLK, and there is no frequency divider inside the module.

\subsection{Wake-Up Modes for Automatic Temperature Monitoring}

There are two wake-up modes for temperature monitoring, selected by \hyperref[fielddesc:TSENSWAKEUPMODE]{TSENS\_WAKEUP\_MODE}:

   \begin{itemize}
       \item Absolute value mode:
           \begin{itemize}
                \item Monitors the absolute value of the current temperature. Configure \hyperref[fielddesc:TSENSWAKEUPTHLOW]{TSENS\_WAKEUP\_TH\_LOW} and \hyperref[fielddesc:TSENSWAKEUPTHHIGH]{TSENS\_WAKEUP\_TH\_HIGH} to set the temperature thresholds. Wake-up will be triggered if the sampled value is above the high threshold or below the low threshold.
           \end{itemize}
       \item Change value mode:
            \begin{itemize}
                \item Monitors the temperature changes inside the chip. If the temperature increment of two consecutive samplings exceeds the high threshold configured in \hyperref[fielddesc:TSENSWAKEUPTHHIGH]{TSENS\_WAKEUP\_TH\_HIGH}, or the temperature decrement of two consecutive samplings exceeds the low threshold configured in \hyperref[fielddesc:TSENSWAKEUPTHLOW]{TSENS\_WAKEUP\_TH\_LOW}, a wake-up will be triggered. For example, when \hyperref[fielddesc:TSENSWAKEUPTHLOW]{TSENS\_WAKEUP\_TH\_LOW} is set to 8, and the two consecutive sampling values are 28 and 19, resulting in a temperature decrement of 9, a wake-up will be triggered.
            \end{itemize}
    \end{itemize}

\subsection{Temperature Measurement Range and Offset}

The temperature sensor is designed with five measurement ranges, as displayed in Table \ref{tab:tsensor-dac}, to enhance accuracy. Each range has its own measurement errors, but users can select particular offsets to obtain calibrated data, so there is no need to be concerned about the errors.

\begin{longtable}[c]{|l|r|}
    \caption{Temperature Measurement Range and Offset}
    \label{tab:tsensor-dac}
    \\\hline
    \rowcolor{lightgray}
    \textbf{Temperature Measurement Range (°C)}   & \textbf{Temperature Offset} \\\hline
    50 \verb+~+ 125 &  –2 \\\hline
    20 \verb+~+ 100 &  –1 \\\hline
    –10 \verb+~+ 80 &  0 \\\hline
    –30 \verb+~+ 50 &  1 \\\hline
    –40 \verb+~+ 20 &  2 \\\hline
\end{longtable}

\subsection{Data Conversion}

The sensor output value (\(VALUE\)) is stored in \hyperref[fielddesc:TSENSOUT]{TSENS\_OUT}. To calculate the actual temperature \(T\) (°C) based on \(VALUE\), use the following equation:
\[
        T = 0.4386 * VALUE – 27.88 * offset – 20.52
\]
where \(offset\) is the temperature offset displayed in Table \ref{tab:tsensor-dac}.

%%%%%%%%%%%%%%%%%%%%%%%%%%%%%%%%%
%%% Event Task Matrix Feature %%%
%%%%%%%%%%%%%%%%%%%%%%%%%%%%%%%%%

\section{Event Task Matrix Feature}
\label{sec:sensor-etm-func}

The temperature sensor on \chipname{} supports the Event Task Matrix (ETM) function, which allows the temperature sensor's ETM tasks to be triggered by any peripherals' ETM events, or temperature sensor's ETM events to trigger any peripherals' ETM tasks. This section introduces the ETM tasks and events related to the temperature sensor. For more information, please refer to Chapter \ref{mod:evntaskmatrix} \nameref{mod:evntaskmatrix}.

The temperature sensor can receive the following ETM tasks:

\begin{itemize}
    \item TMPSNSR\_TASK\_START\_SAMPLE: The temperature sensor starts sampling when this task is triggered.
    \item TMPSNSR\_TASK\_STOP\_SAMPLE: The temperature sensor stops sampling when this task is triggered.
\end{itemize}

The temperature sensor can generate the following ETM events:

\begin{itemize}
    \item TMPSNSR\_EVT\_OVER\_LIMIT: Generated when the temperature is beyond the threshold.
\end{itemize}

In practical applications, the temperature sensor's ETM events can trigger its own ETM tasks. For example, the TMPSNSR\_EVT\_OVER\_LIMIT event can trigger the TMPSNSR\_TASK\_STOP\_SAMPLE task.

%%%%%%%%%%%%%%%%%%%%%
%%% Interrupts %%%
%%%%%%%%%%%%%%%%%%%%%

\section{Interrupts}

\chipname{}'s Temperature Sensor can generate the following interrupt signal that will be sent to the \hyperref[mod:intmtrx]{Interrupt Matrix}.

    \begin{itemize}
        \item LP\_TSENS\_INTR
    \end{itemize}

There is one internal interrupt source from the Temperature Sensor that can generate the above interrupt signal:

    \begin{itemize}
        \item TSENS\_COCPU\_TSENS\_WAKE\_INT\label{int:TSENSINT}: Triggered when the temperature sample value exceeds the threshold.
    \end{itemize}

\begin{tiplisting}
For definitions of \textit{interrupt}, \textit{interrupt signal}, \textit{interrupt source}, and their correlations, please refer to Chapter \ref{mod:intmtrx} \textit{\nameref{mod:intmtrx}} > Section \ref{sec:intmtx-int-terminology} \textit{\nameref{sec:intmtx-int-terminology}}.
\end{tiplisting}

Each interrupt source can be configured by a common set of registers that are described in Section \hyperref[interrupt-config-registers]{\textit{Interrupt Configuration Registers}}. The specific registers can be found in Section \ref{sec:tmpsen-reg-summ} \textit{\nameref{sec:tmpsen-reg-summ}}.

\section{Programming Procedure}

The temperature sensor can be started by software as follows:
    \begin{enumerate}
        \item Set \hyperref[fielddesc:TSENSPOWERUP]{TSENS\_POWER\_UP} to power up the temperature sensor.
        \item Set \hyperref[fielddesc:LPPERICKENLPTSENS]{LPPERI\_CK\_EN\_LP\_TSENS} to enable the temperature sensor clock.
        \item Wait for \hyperref[fielddesc:TSENSXPDWAIT]{TSENS\_XPD\_WAIT} clock cycles till the temperature sensor releases from reset and starts measuring the temperature.
        \item Wait for 100 $\mu$s (so that the output value gradually approaches the actual temperature) and then read the data from \hyperref[fielddesc:TSENSOUT]{TSENS\_OUT}.
    \end{enumerate}

To enable hardware-triggered automatic temperature monitoring, add the following programming steps before powering up the sensor:
    \begin{enumerate}
        \item Configure \hyperref[fielddesc:TSENSSAMPLERATE]{TSENS\_SAMPLE\_RATE} to set sampling rate.
        \item Configure \hyperref[fielddesc:TSENSWAKEUPMODE]{TSENS\_WAKEUP\_MODE} to select wake-up mode for temperature monitoring.
        \item Configure \hyperref[fielddesc:TSENSWAKEUPTHHIGH]{TSENS\_WAKEUP\_TH\_HIGH/LOW} to set the high and low temperature monitoring thresholds.
        \item Set \hyperref[fielddesc:TSENSWAKEUPEN]{TSENS\_WAKEUP\_EN} to start temperature monitoring.
        \item Set \hyperref[fielddesc:TSENSSAMPLEEN]{TSENS\_SAMPLE\_EN} to enable automatic temperature monitoring.
    \end{enumerate}

In automatic temperature monitoring mode, the output value will not be saved. Nevertheless, users can get the value at any moment by accessing \hyperref[fielddesc:TSENSOUT]{TSENS\_OUT}.

% >>> If your module has no registers, delete sample text
%       for Sections Register Summary and Registers below <<<

%%%%%%%%%%%%%%%%%%%%%%%%
%%% Register Summary %%%
%%%%%%%%%%%%%%%%%%%%%%%%

% >>> Update the sample text below and insert
%     the info on register summary into the subfile <<<

\newpage
\hypertarget{tmpsen-reg-summ}{}
\section{Register Summary}
\label{sec:tmpsen-reg-summ}

The addresses in this section are relative to Temperature Sensor base address provided in Table \ref{tab:sysmem-base-address} in Chapter \ref{mod:sysmem} \textit{\nameref{mod:sysmem}}.

The abbreviations given in Column \textbf{Access} are explained in Section \hyperref[glossary-access-types]{\textit{Access Types for Registers}}.

\subfile{\modulefiles/62-TMPSEN-reg-summ__EN}


%%%%%%%%%%%%%%%%%
%%% Registers %%%
%%%%%%%%%%%%%%%%%

% >>> Update the sample text below and insert
%      the info on registers into the subfile <<<

\newpage
\section{Registers}

The addresses in this section are relative to Temperature Sensor base address provided in Table \ref{tab:sysmem-base-address} in Chapter \ref{mod:sysmem} \textit{\nameref{mod:sysmem}}.

For how to program reserved fields, please refer to Section \hyperref[programming-reserved-field]{Programming Reserved Register Field}.

\subfile{\modulefiles/62-TMPSEN-reg__EN}


\end{document}
